% Transcription par lecture visuelle directe, sans OCR.
% Statut : transcrit ; relecture indépendante à effectuer.
% Chapitre transcrit jusqu’au dernier trait de la p. 365, fin du PDF fourni.
% Source : PDF original, vue 174, côté droit, page manuscrite 349 (suite).
\chapter{Fondation de la Congrégation des ouvriers et de celle des dames et rétablissement de la confrérie des Martyrs}

Croissez et multipliez ; tel est l'ordre que la voix puissante du créateur a fait entendre à toutes les créatures après leur avoir donné l'être. Croissez et multipliez ; tel est encore aujourd'hui l'ordre que la même voix fait entendre à la Congrégation ; cette précieuse semence que j'ai remise entre vos mains ne vous a pas été confiée pour vous seuls, mais afin que vous la fassiez fructifier
% Source : PDF original, vue 175, côté gauche, page manuscrite 350.
et rendre au centuple ; j'ai étendu une main tutélaire sur vous, et vous avez vu votre société s'augmenter et se fortifier. Vos premiers membres, parvenus à la maturité de l'âge, sont sortis de vos mains comme un essaim d'abeilles pour former une ruche nouvelle à côté de la ruche mère, et vos forces ont été doublées. Cette Congrégation est semblable à un bel arbre que j'ai planté moi-même et que ma grâce arrose tous les jours ; c'est cette grâce qui lui a fait pousser comme autant de branches vigoureuses toutes vos sections dont les rameaux mettent à couvert sous leurs ombrages les pécheurs, les ignorants, les pauvres, les malades, les prisonniers, les enfants que je leur ai confiés. Remerciez donc ma bonté paternelle de tout le bien qu'elle a opéré parmi vous ; mais je vous le répète, ce n'est point encore assez ; votre mission n'est pas accomplie ; encore une fois, croissez et multipliez ; il faut que toutes les conditions et que tous les sexes entrent dans la joie de leur commun maître. C'est à vous qu'il a été réservé de les y introduire, Congréganistes de 1802 ! Si par une grâce spéciale, j'ai daigné vous choisir pour vous appeler les premiers, c'est afin que fidèles échos, vous fassiez entendre ma voix à d'autres. Souvenez-vous donc de votre antique devise : Allez toujours en avant. Congrégation ! croissez et multipliez.

\section*{Fondation de la Congrégation des ouvriers}
% Mention marginale de la page manuscrite 350 ; « Congrégation des ouvriers » répété p. 352.
C'est ainsi, Messieurs, que Dieu parloit aux cœurs des chefs de notre société. Ils furent attentifs à ces saintes inspirations. En regardant autour d'eux, ils aperçurent bientôt ce qu'il y avoit à faire. Ils se souvinrent que du temps des Jésuites, on voyoit dans notre ville des Congrégations pour tous les genres d'état. Il en existoit même une formée tout exprès, pour ces enfants
% Source : PDF original, vue 175, côté droit, page manuscrite 351.
des villages voisins qui sous le modeste nom d'âniers, venoient chaque jour dans la ville pour en enlever les immondices.

Jusqu'à présent la classe aisée étoit la seule qui pût profiter des avantages offerts aux bons chrétiens par les Congrégations. Mais la classe ouvrière, cette classe si intéressante, où l'on rencontre cependant en grand nombre des hommes remplis de l'esprit de Dieu, ne possédoit point de corps d'élite où en travaillant à sa propre sanctification, on pût essayer de sanctifier aussi les autres. À la vérité il existoit à Lyon sous le nom d'œuvre des hospitaliers, une société bien respectable dont un grand nombre d'ouvriers faisoit partie. Le but de cette société étoit d'exercer des œuvres de charité envers les malades de l'hôtel-Dieu et les prisonniers, en leur rendant des services manuels. Le principal étoit de les peigner et de les raser ; de là cette société étoit connue sous le nom de société des raseurs. Sans doute cet exercice de charité qui après une semaine entière consacrée au travail, devenoit pour de pieux ouvriers, leur seule récréation du dimanche, ne pouvoit qu'être très agréable à celui qui ne laisse pas sans récompense un verre d'eau donné pour l'amour de lui. Mais on ne pouvoit pas y trouver tous ces moyens qu'offrent les Congrégations pour exciter dans les âmes une tendre piété, et un zèle ardent pour la gloire de Dieu. Il n'étoit même que trop vrai que plusieurs membres de cette société réduisoient toute leur religion à leur fidélité à remplir leurs devoirs d'hospitaliers et négligeoient complètement ceux que la loi de Dieu impose à tous les chrétiens.

Cette lacune dans les œuvres de notre ville parut donc à nos préfets devoir appeler toute leur attention. Ils prirent quelques
% Source : PDF original, vue 176, côté gauche, page manuscrite 352.
informations et soit directement, soit indirectement, il leur fut facile de se mettre en rapport avec des ouvriers qui étoient des modèles de piété. On employa le concours de quelques Congréganistes pour faire naître dans leurs cœurs le désir de former une société dans le but de procurer la gloire de Dieu ; cela ne fut pas difficile, car on avoit rencontré des hommes remplis de l'esprit de foi, et disposés à servir la cause de la religion en tout et partout et même au prix des plus grands sacrifices.

Dès qu'on se fut assuré de quelques-uns, on les réunit et le Préfet délégua des commissaires pour représenter auprès d'eux la Congrégation, sans toutefois en compromettre le secret et conférer avec eux des moyens de mener à bonne fin cette grande affaire.

Ces commissaires rédigèrent un projet de règlement, modelé sur les nôtres, sauf les modifications que la différence des états pouvoit exiger. Leurs cœurs étoient si bien préparés que ce fut sans peine qu'on parvint à y introduire toutes les pensées qui avoient présidé à l'établissement de nos Congrégations ; comme nous ils se dévouoient donc de tout leur cœur à la gloire de Dieu, à l'honneur de Marie et à leur propre sanctification. Comme nous, ils considérèrent les bonnes œuvres dont ils pourroient s'occuper, non comme le but spécial de la Congrégation, mais seulement comme un moyen de parvenir à ce but. La Congrégation fut partagée pour l'exercice des bonnes œuvres par divisions comme les nôtres. Sur tout cela, comme sur tous les autres détails de leurs règles, ces braves gens remplis d'humilité ne manifestèrent pas d'autres volontés que celles de nos commissaires. Il n'y eut qu'un seul point auquel ceux-ci
% Source : PDF original, vue 176, côté droit, page manuscrite 353.
rencontrèrent une opposition à laquelle ils ne s'étoient pas attendus ; les chefs de la Congrégation avoient désiré, par des vues peut-être trop humaines, de conserver pour nous seuls notre titre de l'Immaculée Conception ; les commissaires en présentèrent donc un autre, et firent tous leurs efforts pour le faire adopter ; c'étoit le titre de l'Annonciation de la Ste Vierge. On avoit pensé qu'il pouvoit flatter ces braves ouvriers, d'autant mieux que c'est le titre de la première et principale Congrégation établie à Rome qui est considérée comme la mère de toutes les autres. Les commissaires employèrent toute leur rhétorique pour leur persuader d'honorer Marie dans ce grand mystère de l'Incarnation du Verbe, mais tous leurs frais d'éloquence furent complètement perdus. Dieu avoit parlé aux cœurs des futurs Congréganistes. Marie vouloit réunir autour d'elle ses nouveaux enfants sous le même titre qu'elle avoit donné aux anciens. Ils se montrèrent unanimes dans l'intention d'ériger leur Congrégation sous ce seul titre dont nos commissaires ne vouloient pas leur parler, celui de l'Immaculée Conception et à repousser tous les autres. Il fallut bien se soumettre à la volonté de Dieu, si clairement manifestée par la leur.

La fondation de cette nouvelle Congrégation eut lieu, si mes souvenirs sont fidèles, le premier dimanche d'août de l'année 1818. Nos commissaires assistèrent à cette solennité, et furent témoins des engagements des nouveaux Congréganistes. Je ne sais à quoi il faut l'attribuer ; mais malgré tout le zèle que les fondateurs avoient montré à se réunir, leur société eut de la peine à prendre son essor. Elle végéta pendant près de deux années, et demeura presque sans agir pendant tout ce temps.
% Source : PDF original, vue 177, côté gauche, page manuscrite 354.
Cet état de choses ne pouvoit pas durer, car la Providence avoit de grandes vues sur cette Congrégation. La nôtre étoit affligée de voir l'espèce d'engourdissement dans lequel sembloit dormir une société qu'elle considéroit comme sa fille. Mais le ciel avoit placé notre père M. Recorbet comme une sentinelle attentive, chargée de veiller à ce que l'ennemi ne vînt pas apporter des obstacles au succès du bien qui s'entreprenoit sous les auspices de Marie ; mettre enfin en mouvement la Congrégation des ouvriers, étoit certainement une affaire qui méritoit toute son attention. Pour y réussir, il apperçut un moyen dont le succès ne lui paroissoit pas douteux ; c'étoit de la mettre entre les mains d'un directeur actif, capable de réchauffer au besoin, mais surtout de diriger et de féconder le zèle de ces bonnes gens par de sages conseils. Il crut trouver dans M. l'abbé Pater, alors vicaire de St Louis, toutes les qualités nécessaires à ce but désiré ; il soumit sa pensée à M. Courbon, vicaire général, qui s'empressa de désigner M. Pater comme directeur et père spirituel de la Congrégation des ouvriers.

Je ne vous ferai point ici l'éloge de M. Pater. Vous le connoissez, vous savez comme moi, qu'à un zèle éclairé et un esprit de prudence qui ne laissent rien à désirer, il joint le mérite assez rare d'être bon administrateur. Il eût donc été difficile de trouver un meilleur guide pour la nouvelle Congrégation.

Les Préfets déléguèrent en même temps un Congréganiste auprès de cette Congrégation pour l'assister de ses conseils jusqu'à
% Source : PDF original, vue 177, côté droit, page manuscrite 355.
ce qu'une expérience suffisamment prolongée pût lui faire juger qu'elle pouvoit voler de ses propres ailes. Ce délégué a rempli ces fonctions pendant environ deux ans.

Nous aurons occasion de reparler de la Congrégation des ouvriers car nos rapports avec elle se sont beaucoup multipliés dans la suite ; néanmoins laissez-moi vous dire dès à présent que depuis sa fondation, cette société a été constamment pour nous un objet d'édification. On est frappé d'étonnement lorsqu'on considère de près les vertus de la plupart de ses membres. Quant à nous, Messieurs, pardonnez-moi ma franchise, je crois qu'elles doivent aussi nous pénétrer d'une confusion salutaire. En effet en considérant leur zèle, cette abnégation de leur propre volonté, cette humilité sincère avec laquelle ils se placent au-dessous des autres et ont toujours l'air de se compter pour rien. Comment ne pas désirer de voir ces vertus se faire remarquer aussi dans nos Congrégations, et ne pas déplorer que si nous les rencontrons dans quelques-uns de nos amis nous ne puissions pas les admirer dans tous. Réfléchissons sérieusement à cette parole du souverain maître. Malheur aux riches ! et examinons si le désir d'acquérir et le soin de conserver, ne viennent pas trop souvent rétrécir notre cœur et le rendre incapable de ce dévouement complet que nous admirons dans des hommes d'une classe inférieure. Veillons de près afin que ce bagage dont nous sommes accablés ne nous courbe pas vers la terre ; voyons ces vertueux membres de la Congrégation des ouvriers vivant au jour le jour et comptant sur la Providence qui connoît leurs besoins. Ils peuvent courir bien plus vite encore que nous dans la carrière, et cependant elle a été ouverte pour nous avant de l'être pour eux ; efforçons-nous donc de ne pas
% Source : PDF original, vue 178, côté gauche, page manuscrite 356.
demeurer en arrière.

\section*{Congrégation des dames}
% Mention marginale de la page manuscrite 356.
C'est ainsi que la Providence, tout en conduisant la Congrégation dans la voie par laquelle elle doit accomplir la fin pour laquelle elle l'a formée, est toujours attentive à tout ce qui peut procurer notre avancement spirituel. Elle a placé à côté de nous ces braves ouvriers, afin que les exemples touchants de vertu qu'ils nous mettent tous les jours devant les yeux, nous excitent à nous réveiller de notre assoupissement. Elle nous en réserve de non moins utiles dans ceux que nous allons recevoir encore de nos mères, de nos épouses et de nos sœurs, qu'elle veut tout comme nous, appeler à leur tour à travailler à la gloire de notre maître.

Vous vous souvenez sans doute que dans l'histoire de la Congrégation des jeunes gens, nous vous avons parlé d'une Congrégation de dames, fondée par le P. Roger en 1803, et qui quelques années plus tard avoit été transformée en une œuvre de Providence pour l'éducation des jeunes filles.

Les choses en étoient encore au même point, cependant le Seigneur vouloit qu'il existât à Lyon une Congrégation de dames. C'étoit là un complément nécessaire au projet qu'il avoit formé d'y établir une tribu sainte, un peuple choisi, une race royale pour y soutenir le règne de la foi. Il saura donc bien trouver les moyens d'accomplir ses desseins et notre Congrégation n'y sera pas étrangère, car elle a une mission ; il faut qu'elle la remplisse. Une des choses qui m'intéresse le plus en réunissant mes souvenirs pour vous les transmettre, et qui sans doute vous intéressera aussi, c'est de trouver à chaque page que je suis sur le point de transcrire l'occasion d'admirer de nouveau la marche de la Providence, mais rien à cet égard ne me frappe davantage que de la voir souvent se
% Source : PDF original, vue 178, côté droit, page manuscrite 357.
servir des plus minces incidents qui sans leurs résultats eussent passé tout à fait inaperçus, pour amener l'exécution des projets qu'a formés la suprême sagesse. C'est ce que nous allons observer dans les détails relatifs à la fondation de la Congrégation des dames.

Un nom qui a déjà paru dans l'histoire de la Congrégation va se retrouver encore aujourd'hui sous ma plume ; ce nom est celui de notre saint et digne ami, Bertrand du Coin. Je vous ai bien parlé de son courage héroïque et de toutes ses vertus, mais je ne vous ai pas encore tout dit ; il me reste à ajouter que cette auréole de sainteté dont il étoit revêtu, environnoit aussi toute sa famille. Parmi ces âmes privilégiées dans lesquelles Dieu avoit si bien établi le règne de son amour, se trouvoit une sœur de Bertrand du Coin, qui par la générosité de sa foi eût été très volontiers la compagne des souffrances du frère dont elle partageoit le courage et le zèle. C'est à cette vertueuse personne qu'il a été donné d'être la première pierre de l'édifice que le Seigneur a voulu qu'on élevât à sa gloire sous le nom de Congrégation des Dames.

Mademoiselle du Coin appartenoit à cette époque à l'administration de cette œuvre dont nous avons déjà parlé, destinée au soulagement des prisonniers, sous le nom d'œuvre des Charlottes. Elle vint un jour chez un membre de la Congrégation des hommes dans l'unique intention de le consulter sur un fait particulier à cette administration. En ce moment l'esprit saint parla au cœur du Congréganiste, et lui mit les paroles suivantes à la bouche. Au lieu de vous perdre ainsi dans les détails de votre administration des Charlottes, vous feriez bien mieux de chercher à rétablir la Congrégation
% Source : PDF original, vue 179, côté gauche, page manuscrite 358.
des dames, puisque comme vous le savez, l'ancienne a été envahie par l'œuvre de la Providence, alors vous pourriez avoir une section des prisons qui prendroit soin des prisonniers. Mademoiselle du Coin est frappée de cette réflexion, le St Esprit parle aussi à son cœur, et elle répond au Congréganiste sans hésiter un instant : Si votre femme consent à s'en occuper avec moi, j'y suis toute disposée. Le Congréganiste va chercher sa femme, lui fait part des paroles qui viennent d'être échangées ; ces deux dames entrent séance tenante en conférence sur l'affaire qui vient d'être entamée d'une manière si imprévue ; avant d'aller plus loin, elles résolurent cependant de consulter de part et d'autre. Les avis qu'on leur donna furent favorables à l'exécution du projet, en sorte qu'au bout de quelques jours, il ne falloit plus pour pouvoir aller en avant qu'avoir recruté un certain nombre de fondatrices. Ces dames ne tardèrent pas à les rencontrer ; parmi elles se trouvoient les épouses de plusieurs Congréganistes.

Après ces préliminaires qui furent conduits avec assez de rapidité, la fondation de la nouvelle Congrégation put avoir lieu ; je crois que ce fut le 17 septembre 1818. Cette cérémonie se fit dans notre chapelle des Chartreux. La messe fut célébrée par M. Mioland, prêtre Congréganiste alors supérieur des missions du diocèse ; elle fut servie par un membre de la Congrégation des hommes ; la Congrégation s'établit sous le vocable des SS. Cœurs de Jésus et de Marie.

Notre société étoit intervenue dans cette fondation par la première pensée qui avoit été suggérée par un de ses membres. Elle ne se borna pas là. Elle prépara à ces dames
% Source : PDF original, vue 179, côté droit, page manuscrite 359.
un canevas de règlement basé, comme celui des ouvriers, sur le nôtre. M. Mioland fut en premier lieu nommé directeur de la Congrégation naissante. Cette Congrégation eut comme celle des ouvriers, un peu de peine à se mettre en train. Avant de se livrer activement aux bonnes œuvres de tout genre, qu'elle opère maintenant avec tant de succès et surtout avec tant de zèle, elle commença par végéter les premières années. M. Mioland comme vous le savez assez étoit plein de zèle. Étant Congréganiste et connoissant à fond l'esprit qui doit diriger nos sociétés, il eût été certainement bien capable de conduire et de faire prendre son essor à la Congrégation des dames ; mais il étoit trop occupé et trop souvent absent à cause des missions qu'il avoit à donner dans les campagnes, pour qu'il lui fût possible de donner à cette direction la suite nécessaire pour la rendre profitable.

M. Recorbet nous vint encore en aide. M. Pater avoit commencé à diriger la Congrégation des ouvriers et y réussissoit à merveille ; notre père pensa que rien ne l'empêcheroit qu'il pût en même temps diriger la Congrégation des dames. Il engagea donc M. Courbon à l'en charger. M. Mioland s'estima heureux de remettre en des mains aussi sages le dépôt qui lui avoit été confié.

Ces dames songèrent alors pour donner enfin de l'activité à leurs travaux, à mettre la dernière main à leur règlement. Le travail fut préparé par une commission composée du directeur, de quelques-unes des fondatrices et d'un membre de la Congrégation des hommes. Notre société ne négligeoit
% Source : PDF original, vue 180, côté gauche, page manuscrite 360.
pas de saisir toutes les occasions d'inculquer son esprit à toutes les Congrégations qui se formoient autour de nous, car elles avoient toutes le même but et devoient se mouvoir dans la même sphère.

La Congrégation des dames fit dès lors de rapides progrès. Elle s'accrut en nombre. Ses œuvres prirent beaucoup d'extension. Elle s'est constamment montrée animée du meilleur esprit. Nous trouvons encore ici un juste sujet de nous humilier ; car ces dames, malgré la foiblesse de leur sexe et les infirmités auxquelles il est sujet, nous donnent l'exemple d'un zèle et d'une exactitude dans la pratique des bonnes œuvres, qui devroit exciter notre émulation. Que les trésors de mérite que nous voyons amasser sous nos yeux par ceux qui n'ont été appelés qu'après nous à travailler à la vigne du Seigneur réveillent donc enfin notre ardeur. Congréganistes, n'oublions pas, que nous avons été appelés les premiers et ne nous montrons pas indignes de notre belle vocation.

\section*{Confrérie des Martyrs}
% Mention marginale de la page manuscrite 360.
Vers le même temps se forma ou plutôt se rétablit à Lyon, la confrérie des Martyrs. Quoique cette institution ne se rattache pas à la Congrégation d'une manière directe, elle ne lui est pas étrangère ni par son but, ni par les moyens qu'on a employés pour la remettre en exercice. Je ne dois donc pas la passer sous silence ; et d'ailleurs ne sommes-nous pas Lyonnais ? Le culte de nos martyrs n'est-il pas véritablement pour nous le culte des ancêtres ? C'est donc ici un souvenir de famille que nous devons nous estimer heureux de conserver dans nos archives.

Pour honorer ces nombreux martyrs, nos pères dans la foi, il
% Source : PDF original, vue 180, côté droit, page manuscrite 361.
existoit autrefois à Lyon une confrérie célèbre. Cette institution comme toutes celles de ce genre, fut renversée à l'époque de la Révolution. M. Durand en sa qualité de curé de St Irénée étoit préposé à la garde de ces précieuses reliques. Ce saint prêtre, homme rempli de foi et plein de zèle pour la gloire des SS. Martyrs, comprit les obligations que lui imposoient les fonctions qui lui avoient été confiées.

Il conçut la pensée d'appeler les chrétiens héritiers de la foi des Martyrs, à venir de nouveau les entourer de leurs hommages. Pour cela il falloit que le zèle des amis de la religion lui vînt en aide pour travailler avec lui au rétablissement de l'ancienne confrérie. Il commençoit alors d'être d'usage, comme cela se pratique aujourd'hui, dès qu'on avoit besoin d'être secondé pour quelque œuvre de zèle de tourner les yeux du côté de la Congrégation. Le vénérable curé de St Irénée eut donc la pensée de s'y adresser. Dans le désir qui remplissoit son cœur de glorifier nos SS. Martyrs, il lui proposa de devenir elle-même la confrérie des Martyrs.

Une semblable proposition étoit inadmissible. Elle étoit entièrement opposée à notre esprit. Son exécution auroit renversé de fond en comble toute l'économie de notre existence. Elle ne pouvoit donc pas être acceptée. Mais malgré cette impossibilité, l'établissement de la confrérie des Martyrs, pourvu qu'il s'opérât en dehors de la Congrégation, n'en avoit pas moins droit à toutes nos sympathies ; car Marie, notre mère et notre patronne étoit la reine des Martyrs, et les Martyrs eux-mêmes, nous avoient engendrés à la foi au prix de tout le sang qu'ils avoient répandu pour elle. Nos
% Source : PDF original, vue 181, côté gauche, page manuscrite 362.
cœurs émus ne pouvoient donc pas refuser leur concours ; aussi la Congrégation fit-elle tout ce qui étoit en son pouvoir pour seconder les pieuses intentions de M. le Curé. Si elle ne passa pas tout entière avec armes et bagages dans le camp des Martyrs comme il l'auroit désiré, elle s'empressa du moins d'engager tous ses membres à s'enrôler sous leurs drapeaux sans abandonner pour cela celui que nous tenons de notre mère commune. Cette invitation fut accueillie avec tout l'élan qu'elle devoit exciter. Dès les premiers jours soixante Congréganistes se firent inscrire dans la confrérie des Martyrs dont ils composèrent le premier noyau. Cet exemple fut suivi par un grand nombre de bons chrétiens. La confrérie eut de rapides accroissements. Les fidèles de tout sexe, de toute condition, de tous les quartiers de notre ville, considèrent comme un bonheur d'en faire partie.

Des exercices particuliers commencèrent dès lors à avoir lieu tous les mois dans cette vénérable chapelle souterraine, où depuis St Irénée tant de saints ont successivement invoqué le nom du Seigneur ; c'est surtout au jour de la fête des saintes reliques que les pieux confrères, parmi lesquels les Congréganistes se rencontrent toujours en grand nombre, se pressent dans cette crypte si féconde en souvenirs ; c'est dans ce saint lieu que les Martyrs nos ancêtres, laissoient échapper de leurs cœurs enflammés de l'amour de Dieu ces désirs brûlants de procurer sa gloire en se dévouant au Martyre. Leurs soupirs s'élevoient comme les vapeurs parfumées de l'encens, jusqu'au trône de l'Éternel ; ils en faisoient descendre comme une douce rosée ces grâces ineffables qui les remplissoient de force et par
% Source : PDF original, vue 181, côté droit, page manuscrite 363.
lesquelles ils devenoient également insensibles aux séductions comme à la barbarie des Tyrans.

Les Congréganistes viennent s'agenouiller à leur tour à la même place pour implorer aussi les miséricordes de celui qui règne dans tous les siècles. Ils le conjurent le cœur ému et humide, de ne pas permettre que dans ces jours de défection ils abandonnent jamais la cause de la foi. Oui, mon Dieu ! s'écrient-ils, étendez sur votre Église une main secourable, défendez-la contre tous ses ennemis ; soutenez en particulier cette cité des Martyrs ; que cette ville qui vous offrit jadis en holocauste le sang d'un si grand nombre de ses enfants ne perde jamais ce précieux trésor de la foi qu'ils lui ont laissé pour héritage. Protégez surtout, ô mon Dieu, cette petite Congrégation que vous avez formée vous-même dans votre miséricorde, et qui vient en ce lieu se confondre avec tous les pieux confrères des Martyrs ; qu'elle se montre toujours digne de votre mère qui a bien voulu être la sienne ; que sous ses auspices tous ceux qui la composent trouvent toujours la force nécessaire pour triompher des obstacles qui pourront s'opposer à leur salut.

Vous l'avez éprouvé tout comme moi, Messieurs, lorsqu'on se trouve dans cette sainte assemblée, on se croit de nouveau transporté au second siècle de l'Église, quand l'évêque Irénée offroit au même lieu les saints mystères, en présence de son peuple. Tout comme alors, les fidèles viennent encore aujourd'hui, par la participation commune au saint banquet, se renouveler tous ensemble dans le sang de l'Agneau, arrosée de nouveau par ce sang précieux qui fait germer
% Source : PDF original, vue 182, côté gauche, page manuscrite 364.
toutes les vertus dans les cœurs, la foi se retrempe et retrouve sa première vigueur. Cet arbre vivra, ainsi que le saint évêque d'Alger le disoit naguère dans ce lieu même, cet arbre qui a jeté sur notre sol de si profondes racines, vivra toujours dans notre patrie. Nous acceptons avec une vive confiance l'assurance que ce vertueux pontife nous en a donnée. Oui, jusqu'à la fin des siècles, cet arbre chéri produira parmi nos descendants des fleurs et des fruits. Puisse la Congrégation ne demeurer jamais étrangère à l'accomplissement de cette prophétie.

Cette solennité est tout intérieure. Elle donne chaque année aux confrères l'occasion de s'édifier entre eux par les exemples de piété qu'ils reçoivent les uns des autres ; elle est pour les Congréganistes comme une fête de plus ajoutée à celles de la Congrégation avec lesquelles elle a beaucoup d'analogie. Mais il en est une autre dans laquelle Confrères et Congréganistes, car ici ils ne font qu'un, ne sont pas seulement appelés à s'édifier entre eux, mais aussi à édifier leurs concitoyens et à répandre la bonne odeur de J. C. par leur assistance à une procession solennelle.

L'étendard des Martyrs est déployé ; on y lit en lettres d'or cette sentence confirmée par l'histoire de l'Église de tous les siècles : \textit{Sanguis Martyrum, semen Christianorum}\footnote{Littéralement : « Le sang des martyrs, semence des chrétiens ». Source et forme précise de la sentence à vérifier, Q253. (Note du transcripteur).}. Les confrères suivent avec joie ce glorieux trophée qu'ils contemplent avec émotion. Au milieu de leurs rangs de pieux lévites portent en triomphe les précieuses reliques de notre apôtre, de notre père St Irénée. Cette pompe majestueuse parcourt les rues de l'antique Lugdunum, de cette cité sainte, empourprée du sang des Martyrs. Ces chrétiens
% Source : PDF original, vue 182, côté droit, page manuscrite 365.
du 19\up{e} siècle font retentir les airs des invocations qu'ils adressent aux héros du second siècle, à ces saints qui ont habité ainsi que nous habitons, qui ont respiré l'air que nous respirons, qui ont logé dans les maisons où nous logeons, et enfin qui ont professé et qui sont morts pour la foi que nous professons. On croit encore apercevoir dans ce sol l'empreinte des pas des Pothin, des Irénée, des Blandine, des Attale, des Épipode, des Alexandre et de tant d'autres illustres serviteurs de notre Dieu. Il semble que cette terre ainsi sanctifiée, doit en ce moment rendre saints les enfants des saints.

Les cœurs des Congréganistes raniment leur foi en s'unissant au siècle des Martyrs, mais en même temps ils s'élancent avec amour jusqu'aux pieds de Marie, leur mère, qui fut toujours celle de nos pères dans la foi. Ô Marie ! s'écrient-ils dans un saint transport, notre puissante avocate, notre secours, notre refuge, entendez nos prières et nos vœux ; unissez-les aux mérites de [mot incertain : vos] ancêtres\footnote{Le possessif proposé « vos » reste à confirmer, Q254 ; pas de substitution d'après le sens attendu. (Note du transcripteur).}, d'amples récompenses les attendent dans le séjour de la gloire. Que les exemples qu'ils nous ont laissés nous apprennent à rester fidèles jusqu'à la mort, et à ne jamais cesser de tout rapporter à celui qui est admirable dans ses saints et le seul saint par essence. Obtenez-nous aussi la grâce de ne jamais oublier que si nous persévérons jusqu'à la fin, nous partagerons un jour le bonheur dont nos glorieux Martyrs sont en possession, et alors comme eux et avec eux nous verrons Dieu, nous l'aimerons, nous le bénirons dans les siècles des siècles et dans l'éternité tout entière. Nous chanterons aussi vos louanges, ô Marie, car c'est à vous que nous devons tout notre bonheur.

% Un trait horizontal clôt le texte au bas de la page 365, dernière page du PDF fourni.
